% ---------------------------------------------------------------------
% 2.2 Electroencefalografía
% ---------------------------------------------------------------------
\section[ELECTROENCEFALOGRAFÍA]{Electroencefalografía}
\label{sec:eeg}

El electroencefalograma (EEG) se define como el registro de la actividad eléctrica del cerebro
tomado por medio de electrodos colocados en la superficie de la cabeza, mientras que la
electroencefalografía corresponde a la disciplina que se encarga de registrar e interpretar dichos
registros \cite{kane2017}. El primer registro de este tipo en humanos fue reportado en 1929 por el
psiquiatra alemán Hans Berger, quien además acuñó el término electroencefalograma y describió el
ritmo que hoy se conoce como alfa \cite{berger1929}.

\subsection{Adquisición de la señal}

Los potenciales registrados en el cuero cabelludo tienen amplitudes que van de unos cuantos
microvolts hasta alrededor de 100~$\mu$V, por lo que deben amplificarse considerablemente antes
de digitalizarse \cite{sanei2007}. Los electrodos más comunes son discos de plata recubiertos de
cloruro de plata (Ag/AgCl). Se colocan con un gel conductor para reducir la impedancia de contacto,
que idealmente debe ser menor a 5~k$\Omega$. También existen electrodos secos, que no necesitan gel
y se colocan más rápido. Su desventaja es que tienen mayor impedancia y son más sensibles a los
artefactos de movimiento \cite{niedermeyer2005}. En este sentido, el trabajo de Contreras \cite{contreras2021} utilizó
electrodos secos de un sistema de bajo costo con cuatro canales, mientras que las bases de datos
públicas del paradigma P300, como la empleada en este trabajo, fueron adquiridas con equipos de
grado de investigación de 64 canales \cite{blankertz2006}.

Cada canal de EEG mide en realidad una diferencia de potencial, por lo que es necesario definir
contra qué punto se mide. En un montaje unipolar, cada electrodo se compara con un electrodo
de referencia común, ubicado en un sitio con poca actividad cerebral, como el lóbulo de la oreja o
las mastoides. En un montaje bipolar, en cambio, se registra la diferencia entre pares de
electrodos vecinos \cite{ganong2010}. Una tercera
alternativa, muy utilizada en el procesamiento digital, es la referencia promedio común (CAR), en
la que a cada canal se le resta el promedio de todos los canales en cada instante de tiempo
\cite{luck2014}.

\subsection{Sistema internacional 10-20}

Para que los registros sean comparables entre sujetos y laboratorios, la posición de los
electrodos se estandariza mediante el sistema internacional 10-20, propuesto originalmente por
Jasper \cite{jasper1958} y actualizado por la Federación Internacional de Neurofisiología Clínica
(IFCN) \cite{seeck2017}. Su nombre se debe a que las posiciones se obtienen dividiendo las
distancias entre cuatro puntos de referencia anatómicos (el nasión, el inión y los dos puntos
preauriculares) en segmentos del 10\% y del 20\% de su longitud total. Cada electrodo se nombra con
una o dos letras según la región que cubre: Fp (frontopolar), F (frontal), C (central), T
(temporal), P (parietal) y O (occipital). Después se agrega un número impar si el electrodo está en
el hemisferio izquierdo, un número par si está en el derecho, o la letra \textit{z} si está sobre
la línea media \cite{seeck2017}. La Figura~\ref{fig:1020} muestra la
distribución de los 21 electrodos del sistema, resaltando los de la línea media (Fz, Cz, Pz y Oz),
que fueron los utilizados en el trabajo de Contreras \cite{contreras2021}. Las bases de datos del
paradigma P300 suelen utilizar la extensión 10-10, que agrega posiciones intermedias hasta
completar 64 o más canales \cite{blankertz2006}.

\begin{figure}[htbp]
\centering
\includegraphics[width=0.52\textwidth]{02-marco-referencial/sistema_10_20.pdf}
\caption{Posiciones de los electrodos del sistema internacional 10-20 (vista superior). Se
resaltan los electrodos de la línea media Fz, Cz, Pz y Oz. Elaboración propia con base en
\cite{seeck2017}.}
\label{fig:1020}
\end{figure}

\subsection{Ritmos cerebrales}

Un ritmo cerebral se define como la actividad del EEG formada por ondas de un periodo
aproximadamente constante \cite{kane2017}, y de acuerdo con su frecuencia se clasifica en las
bandas delta, theta, alfa, beta y gamma. Cada banda se asocia con estados de actividad distintos:
el ritmo alfa, por ejemplo, predomina en la región occipital cuando la persona está despierta, en
reposo y con los ojos cerrados, y se sustituye por actividad beta de menor amplitud cuando la
atención se dirige a una tarea, fenómeno conocido como bloqueo alfa \cite{guyton2016,ganong2010}.
La Tabla~\ref{tab:ritmos} resume sus
características y la Figura~\ref{fig:ritmos} ilustra su forma de onda.

\begin{table}[htbp]
\centering
\caption{Bandas de frecuencia del EEG}
\label{tab:ritmos}
\small
\begin{tabularx}{\textwidth}{L{1.7cm}L{1.85cm}L{1.9cm}L{2.55cm}Y}
\toprule
\textbf{Ritmo} & \textbf{Frecuencia} & \textbf{Amplitud} & \textbf{Región} & \textbf{Estado asociado} \\
\midrule
Delta ($\delta$) & 0.5--4 Hz & 20--200 $\mu$V & Variable & Sueño profundo, lactancia; en vigilia puede indicar patología. \\
Theta ($\theta$) & 4--8 Hz & 10--50 $\mu$V & Parietal, temporal & Somnolencia, niñez, algunos estados de estrés emocional. \\
Alfa ($\alpha$) & 8--13 Hz & 50--100 $\mu$V & Occipital & Vigilia relajada con ojos cerrados; se bloquea al atender. \\
Beta ($\beta$) & 13--30 Hz & $<$30 $\mu$V & Frontal, central & Alerta, concentración y actividad mental. \\
Gamma ($\gamma$) & $>$30 Hz & $<$10 $\mu$V & Difusa & Procesamiento sensorial y atención focalizada. \\
\bottomrule
\end{tabularx}
\par\smallskip{\footnotesize Elaboración propia con información de \cite{guyton2016,ganong2010,kane2017}.}
\end{table}

\begin{figure}[htbp]
\centering
\includegraphics[width=0.85\textwidth]{02-marco-referencial/ritmos_eeg.pdf}
\caption{Forma de onda ilustrativa de los principales ritmos del EEG. Elaboración propia.}
\label{fig:ritmos}
\end{figure}

\subsection{Artefactos en la señal de EEG}

Debido a su baja amplitud, el EEG es muy susceptible a señales que no provienen de la actividad
cerebral, conocidas como artefactos. Los más frecuentes son de origen fisiológico, como los
parpadeos y movimientos oculares, cuya amplitud puede ser varias veces mayor que la del propio
EEG, la actividad eléctrica de los músculos de la cara y el cuello o la actividad cardiaca. Otros que se pueden
presentar son los de origen técnico, como la interferencia de la red eléctrica, que en México se presenta a
60~Hz, o los cambios de impedancia de los electrodos \cite{sanei2007,niedermeyer2005}. La
Tabla~\ref{tab:artefactos} presenta los principales artefactos junto con las técnicas más comunes
para mitigarlos, algunas de las cuales se retoman en la Sección~\ref{sec:pds}.

\begin{table}[htbp]
\centering
\caption{Artefactos comunes en el registro de EEG}
\label{tab:artefactos}
\small
\begin{tabularx}{\textwidth}{L{2.8cm}L{2.3cm}L{2.4cm}Y}
\toprule
\textbf{Artefacto} & \textbf{Origen} & \textbf{Banda afectada} & \textbf{Mitigación} \\
\midrule
Parpadeo y movimiento ocular (EOG) & Fisiológico & $<$4 Hz & Filtrado pasa altas, regresión con canales EOG, ICA. \\
Actividad muscular (EMG) & Fisiológico & $>$20 Hz & Filtrado pasa bajas, instrucciones al sujeto, ICA. \\
Actividad cardiaca (ECG) & Fisiológico & 1--2 Hz y armónicos & Referencia adecuada, ICA. \\
Sudoración y deriva de línea base & Fisiológico / electrodo & $<$0.5 Hz & Filtrado pasa altas, corrección de línea base. \\
Interferencia de la red eléctrica & Técnico & 60 Hz (México) & Filtro notch, blindaje, buena impedancia. \\
Movimiento o mal contacto del electrodo & Técnico & Banda ancha & Rechazo de épocas por umbral de amplitud. \\
\bottomrule
\end{tabularx}
\par\smallskip{\footnotesize Elaboración propia con información de \cite{sanei2007,niedermeyer2005,jung2000}.}
\end{table}
